\part{Alg\`ebre et Algorithmique}

%% ════════════════════════════════════════════════════════════════
\chapter{Logique}
\begin{prerequis}
Savoir distinguer une phrase affirmative d'une question; connaître les symboles
$=$, $\ne$, $<$, $\leq$; savoir produire un contre-exemple simple.
\end{prerequis}
\begin{automatismes}
Dire si chaque phrase est vraie, fausse, ou si elle ne possède pas de valeur de
vérité : « $7$ est premier »; « $x+2=5$ »; « calculer $3^2$ »; « tout carré est
positif »; « il existe un entier pair supérieur à $100$ ».
\end{automatismes}
\begin{activite}[Activité d'entrée -- Le tribunal des affirmations]
Par groupes, classer dix phrases proposées par le professeur. Pour chaque phrase
classée vraie ou fausse, rédiger une justification; pour chaque phrase refusée,
expliquer quelle information manque. Transformer ensuite deux phrases ouvertes
en propositions à l'aide de « pour tout » ou « il existe ».
\end{activite}


%% ════════════════════════════════════════════════════════════════

\begin{anecdote}[Histoire des math\'ematiques]
Aristote ($384$--$322$ av.\ J.-C.) fonda la logique formelle. Deux mill\'enaires plus
tard, George Boole ($1815$--$1864$) eut l'id\'ee de traiter la logique comme une
alg\`ebre --- naissance de l'alg\`ebre de Boole, fondement de tous les ordinateurs
modernes. Chaque fois que vous allumez un smartphone, vous utilisez la logique
math\'ematique.
\end{anecdote}

\begin{objectifs}
\begin{itemize}
  \item Distinguer \'enonc\'e, proposition et pr\'edicat.
  \item Ma\^itriser les connecteurs logiques et leurs tables de v\'erit\'e.
  \item Utiliser les quantificateurs $\forall$ et $\exists$.
  \item Former la n\'egation et la r\'eciproque d'une implication.
  \item D\'ecouvrir, en approfondissement, la contrapos\'ee et le raisonnement par l'absurde.
\end{itemize}
\end{objectifs}

\section{\'Enonc\'es et propositions}

\begin{defn}[Proposition]
Une \textbf{proposition} est un \'enonc\'e pouvant \^etre \textbf{vrai}~(\vrai)
ou \textbf{faux}~(\faux), mais pas les deux simultan\'ement.
\end{defn}

\begin{ex}
\begin{itemize}
  \item \og\(2 + 2 = 4\)\fg{} : proposition vraie.
  \item \og\(7\) est pair\fg{} : proposition fausse.
  \item \og\(x > 3\)\fg{} : pr\'edicat (sa valeur d\'epend de $x$), \emph{pas} une proposition.
  \item \og Viens ici\fg{} : phrase imp\'erative, pas une proposition.
\end{itemize}
\end{ex}

\section{N\'egation et connecteurs logiques}

\begin{defn}[Connecteurs logiques]
Soient $P$ et $Q$ deux propositions. On d\'efinit :

\medskip\centering
\renewcommand{\arraystretch}{1.4}
\begin{tabular}{|c|c|c|c|c|c|}
  \hline
  \rowcolor{BN!15}
  $P$ & $Q$ & $\lnot P$ & $P \land Q$ & $P \lor Q$ & $P \Rightarrow Q$ \\
  \hline
  \vrai & \vrai & \faux & \vrai & \vrai & \vrai \\
  \vrai & \faux & \faux & \faux & \vrai & \faux \\
  \faux & \vrai & \vrai & \faux & \vrai & \vrai \\
  \faux & \faux & \vrai & \faux & \faux & \vrai \\
  \hline
\end{tabular}

\medskip\flushleft
L'implication $P \Rightarrow Q$ est fausse \emph{uniquement} quand $P$~est vraie et
$Q$~est fausse. Le \og ou\fg\ math\'ematique est \textbf{inclusif}.
\end{defn}

\begin{defn}[\'Equivalence]
$P \Leftrightarrow Q$ est vraie si $P$ et $Q$ ont la m\^eme valeur de v\'erit\'e :
$$P \Leftrightarrow Q \;\equiv\; (P \Rightarrow Q) \land (Q \Rightarrow P).$$
\end{defn}

\horsprogramme{La contrapos\'ee et le raisonnement par l'absurde sont propos\'es ici comme enrichissement. Ils ne doivent pas masquer les raisonnements directs attendus en T10.}

\begin{methode}[Contrapos\'ee]
La \textbf{contrapos\'ee} de $P \Rightarrow Q$ est $\lnot Q \Rightarrow \lnot P$.
Ces deux implications sont toujours \'equivalentes. Pour d\'emontrer
$P \Rightarrow Q$, il suffit de d\'emontrer $\lnot Q \Rightarrow \lnot P$.
\end{methode}

\begin{figure}[H]
\centering
\begin{tikzpicture}[
  node distance=0.8cm and 1.4cm,
  decision/.style={draw=OR,fill=OR!10,diamond,aspect=2.5,
                   minimum width=3.5cm,minimum height=0.9cm,
                   align=center,font=\small},
  action/.style={draw=BN,fill=BN!8,rounded corners=3pt,
                 minimum width=3.3cm,minimum height=0.75cm,
                 align=center,font=\small},
  result/.style={draw=VE,fill=VE!10,rounded corners=4pt,
                 minimum width=3.3cm,minimum height=0.75cm,
                 align=center,font=\small\bfseries},
  arr/.style={-Stealth,BN!70,thick}]

  \node[action,fill=BN,text=white] (start) {D\'emontrer $P\Rightarrow Q$};
  \node[decision,below=of start] (q1) {$\lnot Q\Rightarrow\lnot P$\\facile~?};
  \node[result,below right=0.8cm and 1.6cm of q1] (contra)
    {Contrapos\'ee\\$\lnot Q\Rightarrow\lnot P$};
  \node[decision,below=1.6cm of q1] (q2) {Supposer $P\land\lnot Q$\\m\`ene \`a une\\contradiction~?};
  \node[result,below right=0.8cm and 1.6cm of q2] (absurde)
    {Absurde~:\\supposer $P\land\lnot Q$,\\trouver $\bot$};
  \node[result,below=1.6cm of q2] (direct)
    {Direct~:\\supposer $P$,\\d\'eduire $Q$};

  \draw[arr] (start)--(q1);
  \draw[arr] (q1) -- node[above right,font=\scriptsize]{Oui} (contra);
  \draw[arr] (q1) -- node[left,font=\scriptsize]{Non} (q2);
  \draw[arr] (q2) -- node[above right,font=\scriptsize]{Oui} (absurde);
  \draw[arr] (q2) -- node[left,font=\scriptsize]{Non} (direct);
\end{tikzpicture}
\caption{Choisir sa m\'ethode de d\'emonstration.
Commencer par chercher si la contrapos\'ee est plus simple~;
si non, tenter l'absurde~; sinon, partir directement de $P$.}
\end{figure}

\begin{attention}
Ne pas confondre une implication et sa r\'eciproque.
\og\(x=2 \Rightarrow x^2=4\)\fg{} est vraie.
\og\(x^2=4 \Rightarrow x=2\)\fg{} est \textbf{fausse} (contre-exemple : $x=-2$).
\end{attention}

\section{Quantificateurs}

\begin{defn}[Quantificateurs $\forall$ et $\exists$]
\begin{itemize}
  \item $\forall x \in E,\ P(x)$ : \og pour tout $x$ dans $E$, $P(x)$ est vraie.\fg
  \item $\exists x \in E,\ P(x)$ : \og il existe au moins un $x$ dans $E$ tel que
        $P(x)$ est vraie.\fg
\end{itemize}
\end{defn}

\begin{prop}[N\'egations des quantificateurs]
$$\lnot\bigl(\forall x \in E,\ P(x)\bigr) \;\equiv\; \exists x \in E,\ \lnot P(x)$$
$$\lnot\bigl(\exists x \in E,\ P(x)\bigr) \;\equiv\; \forall x \in E,\ \lnot P(x)$$
\end{prop}

% Diagramme de Venn
\begin{figure}[H]
\centering
\begin{tikzpicture}[scale=1.05]
  \draw[BN,thick] (-3.8,-2.1) rectangle (3.8,2.1);
  \node[BN,font=\small] at (3.4,1.8) {$\Omega$};
  \fill[TQ!25] (-0.9,0) circle (1.6cm);
  \fill[OR!25] (0.9,0) circle (1.6cm);
  \begin{scope}
    \clip (-0.9,0) circle (1.6cm);
    \fill[TQ!35!OR!35] (0.9,0) circle (1.6cm);
  \end{scope}
  \draw[TQ,thick] (-0.9,0) circle (1.6cm);
  \draw[OR,thick] (0.9,0) circle (1.6cm);
  \node[font=\bfseries,text=TQ] at (-2,0) {$P$};
  \node[font=\bfseries,text=OR] at (2,0) {$Q$};
  \node[font=\small,text=BN] at (0,0) {$P \land Q$};
  \node[font=\footnotesize,text=GM!70!black,align=center] at (0,-2.7)
    {L'union des deux cercles est $P \lor Q$.\quad
     L'intersection est $P \land Q$.};
\end{tikzpicture}
\caption{Diagramme de Venn : conjonction et disjonction.}
\end{figure}

\begin{figure}[H]
\centering
\renewcommand{\arraystretch}{1.5}
\begin{tabular}{|>{\columncolor{BN!8}}c|>{\columncolor{BN!8}}c||c|c|c|c|}
\hline
\rowcolor{BN!18}
$P$ & $Q$ & $\lnot P$ & $P\land Q$ & $P\lor Q$ & $P\Rightarrow Q$ \\
\hline
\textcolor{VE!80!black}{\textbf{V}} & \textcolor{VE!80!black}{\textbf{V}}
  & \textcolor{BO}{\textbf{F}}
  & \textcolor{VE!80!black}{\textbf{V}}
  & \textcolor{VE!80!black}{\textbf{V}}
  & \textcolor{VE!80!black}{\textbf{V}} \\
\textcolor{VE!80!black}{\textbf{V}} & \textcolor{BO}{\textbf{F}}
  & \textcolor{BO}{\textbf{F}}
  & \textcolor{BO}{\textbf{F}}
  & \textcolor{VE!80!black}{\textbf{V}}
  & \textcolor{BO}{\textbf{F}} \\
\rowcolor{BN!4}
\textcolor{BO}{\textbf{F}} & \textcolor{VE!80!black}{\textbf{V}}
  & \textcolor{VE!80!black}{\textbf{V}}
  & \textcolor{BO}{\textbf{F}}
  & \textcolor{VE!80!black}{\textbf{V}}
  & \textcolor{VE!80!black}{\textbf{V}} \\
\rowcolor{BN!4}
\textcolor{BO}{\textbf{F}} & \textcolor{BO}{\textbf{F}}
  & \textcolor{VE!80!black}{\textbf{V}}
  & \textcolor{BO}{\textbf{F}}
  & \textcolor{BO}{\textbf{F}}
  & \textcolor{VE!80!black}{\textbf{V}} \\
\hline
\end{tabular}
\caption{Tables de v\'erit\'e des connecteurs principaux.
\textcolor{VE!80!black}{\textbf{V}} = vrai~; \textcolor{BO}{\textbf{F}} = faux.
$P\Rightarrow Q$ est fausse \emph{uniquement} quand $P$ est vraie et $Q$ est fausse.}
\end{figure}

%─────────────────────────────────────────────────────────────────
\horsprogramme{Ces schémas de preuve permettent de consolider le raisonnement direct et de découvrir d'autres démarches de démonstration.}
\section{Types de raisonnement}
%─────────────────────────────────────────────────────────────────

\begin{methode}[Raisonnement direct]
Pour d\'emontrer $P \Rightarrow Q$ directement :
on \textbf{suppose $P$ vraie} et on en d\'eduit, par une cha\^ine d'implications,
que $Q$ est vraie.
\end{methode}

\begin{methode}[Raisonnement par contrapos\'ee]
On d\'emontre $\lnot Q \Rightarrow \lnot P$ \`a la place de $P \Rightarrow Q$.
\textbf{\'Etapes :} (1) Supposer $\lnot Q$. (2) En d\'eduire $\lnot P$.
\end{methode}

\begin{methode}[Raisonnement par l'absurde]
Pour d\'emontrer $P$ :
on \textbf{suppose $\lnot P$} et on aboutit \`a une \textbf{contradiction}.
La contradiction montre que $\lnot P$ est fausse, donc $P$ est vraie.
\end{methode}

\begin{attention}
Un contre-exemple suffit pour r\'efuter $\forall x \in E,\ P(x)$.
En revanche, un seul exemple ne suffit pas pour prouver une propri\'et\'e
universelle.
\end{attention}

\begin{ex}[Raisonnement par l'absurde]
D\'emontrons que $\sqrt{2}$ est irrationnel.
\emph{Supposons} $\sqrt{2} \in \Q$ : il existe $p, q \in \N^*$ avec $\pgcd(p,q)=1$
et $\sqrt{2} = p/q$.
Alors $2q^2 = p^2$, donc $p^2$ est pair, donc $p$ est pair : $p = 2k$.
On obtient $2q^2 = 4k^2$, d'o\`u $q^2 = 2k^2$, et $q$ est pair.
Mais alors $2 \mid \pgcd(p,q)$, ce qui contredit $\pgcd(p,q)=1$. \textbf{Contradiction.}
Donc $\sqrt{2} \notin \Q$.
\end{ex}

%─────────────────────────────────────────────────────────────────
\section{Exercices}
%─────────────────────────────────────────────────────────────────

\subsection*{\diff{1}\quad Niveau 1 \textemdash\ Bases et reconnaissance}

\begin{exo}[\logiq\ Propositions ou non \corrige]{1}
Parmi les \'enonc\'es ci-dessous, indiquer lesquels sont des \textbf{propositions}.
Si c'est une proposition, pr\'eciser sa valeur de v\'erit\'e.
\index{proposition}
\begin{multicols}{2}
\begin{enumerate}
  \item $5 \times 4 = 20$
  \item $\sqrt{2}$ est un entier
  \item $x + 3 = 7$
  \item Tout triangle \'equilat\'eral est isoc\`ele
  \item $0 \times 999 = 0$
  \item Bonjour, comment vas-tu~?
  \item $\pi > 3$
  \item $n$ est premier
  \item $\forall n \in \N,\ n \geq 0$
  \item Ferme la porte.
\end{enumerate}
\end{multicols}
\end{exo}

\begin{exo}[\logiq\ N\'egations simples \corrige]{1}
Donner la n\'egation\index{négation} de chaque proposition.
\begin{enumerate}
  \item $17$ est un nombre premier.
  \item Le triangle $ABC$ est rectangle.
  \item $x \geq 5$.
  \item Il pleut \textbf{et} il fait froid.
  \item Il pleut \textbf{ou} il fait chaud.
  \item $\forall n \in \N,\ n^2 \geq n$.
  \item $\exists x \in \R,\ x^2 = -1$.
  \item $2n$ est pair pour tout entier $n$.
\end{enumerate}
\end{exo}

\begin{exo}[\logiq\ Tables de v\'erit\'e \`a deux variables \corrige]{1}
Dresser la table de v\'erit\'e compl\`ete des formules suivantes.
\begin{multicols}{2}
\begin{enumerate}
  \item $\lnot P \lor Q$
  \item $P \land \lnot Q$
  \item $(P \lor Q) \land \lnot P$
  \item $\lnot(P \land Q)$
  \item $\lnot P \land \lnot Q$
  \item $\lnot P \lor \lnot Q$
\end{enumerate}
\end{multicols}
Comparer les formules (4) et (6). Quelle est la n\'egation de
$(P\lor Q)$~? (Lois de De Morgan\index{lois de De Morgan}.)
\end{exo}

\begin{exo}[Valeur de v\'erit\'e d'une implication \corrige]{1}
Chacune des implications suivantes est-elle vraie ou fausse~?
Justifier bri\`evement.
\begin{enumerate}
  \item Si $2 + 2 = 4$, alors $3 + 3 = 6$.
  \item Si $2 + 2 = 5$, alors la Terre est plate.
  \item Si $2 + 2 = 4$, alors $1 = 0$.
  \item Si $n$ est divisible par $4$, alors $n$ est divisible par $2$.
  \item Si $n$ est divisible par $2$, alors $n$ est divisible par $4$.
\end{enumerate}
\end{exo}

\begin{exo}[\logiq\ Quantificateurs : traduction \corrige]{1}
Traduire en fran\c cais, puis d\'eterminer la valeur de v\'erit\'e.
\begin{enumerate}
  \item $\forall x \in \N,\ x+1 > x$
  \item $\exists n \in \Z,\ n^2 = -4$
  \item $\forall x \in \R,\ x^2 \geq 0$
  \item $\exists x \in \Q,\ x^2 = 2$
  \item $\forall x \in \R^*,\ 1/x > 0$
  \item $\exists n \in \N,\ n^2 = 2n$
\end{enumerate}
\end{exo}

\begin{exo}[\'Ecrire avec des quantificateurs \corrige]{1}
Traduire les propositions suivantes en utilisant $\forall$ et $\exists$.
\begin{enumerate}
  \item Tout carr\'e est un rectangle.
  \item Il existe un entier pair qui est premier.
  \item La somme de deux entiers pairs est toujours paire.
  \item Tout nombre r\'eel admet un oppos\'e.
  \item Il n'existe pas d'entier strictement compris entre $0$ et $1$.
\end{enumerate}
\end{exo}

\begin{exo}[N\'egation avec quantificateurs \corrige]{1}
\'Ecrire la n\'egation de chaque proposition, puis d\'eterminer la valeur
de v\'erit\'e de la n\'egation.
\begin{enumerate}
  \item $\forall x \in \R,\ x^2 > 0$
  \item $\exists n \in \N,\ 2n+1 = 0$
  \item $\forall n \in \N,\ n^2 \geq n$
  \item $\exists x \in \R,\ x^3 = -8$
  \item $\forall x \in \R,\; \forall y \in \R,\ x + y = y + x$
\end{enumerate}
\end{exo}

\begin{exo}[\logiq\ R\'eciproque et contrapos\'ee \corrige]{1}
Pour chaque implication, \'ecrire (a)~la r\'eciproque\index{réciproque},
(b)~la contrapos\'ee\index{contraposée}.
Indiquer laquelle est vraie parmi l'implication et sa r\'eciproque.
\begin{enumerate}
  \item Si $ABC$ est \'equilat\'eral, alors $ABC$ est isoc\`ele.
  \item Si $x = 3$, alors $x^2 = 9$.
  \item Si $n$ est divisible par $6$, alors $n$ est divisible par $3$.
  \item Si $x > 1$, alors $x^2 > 1$.
\end{enumerate}
\end{exo}

\begin{exo}[Contre-exemples \corrige]{1}
R\'efuter chacune des propositions par un contre-exemple\index{contre-exemple}.
\begin{enumerate}
  \item $\forall n \in \N,\ n^2 > n$
  \item $\forall x \in \R,\ \sqrt{x^2} = x$
  \item Si $x^2 = 4$, alors $x = 2$.
  \item Tout multiple de $4$ est un multiple de $8$.
  \item $\forall n \in \Z,\ |n| = n$
\end{enumerate}
\end{exo}

\begin{exo}[Implication ou \'equivalence~? \corrige]{1}
D\'eterminer si l'on a $(\Rightarrow)$, $(\Leftarrow)$ ou $(\Leftrightarrow)$.
Justifier par une d\'emonstration ou un contre-exemple.
\begin{enumerate}
  \item $x = 2 \quad \text{et} \quad x^2 = 4$
  \item $n$ pair \quad et \quad $n^2$ pair
  \item $n$ divisible par $6$ \quad et \quad $n$ divisible par $2$
  \item $|x| = 3$ \quad et \quad $x = 3$
  \item $xy = 0$ \quad et \quad $x = 0$
\end{enumerate}
\end{exo}

\begin{exo}[Table : contrapos\'ee \corrige]{1}
V\'erifier par table de v\'erit\'e que $P \Rightarrow Q$ et
$\lnot Q \Rightarrow \lnot P$ ont toujours la m\^eme valeur de v\'erit\'e.
\end{exo}

\begin{exo}[Table : connecteur $\Leftrightarrow$ \corrige]{1}
V\'erifier par table de v\'erit\'e :
$$P \Leftrightarrow Q \;\equiv\; (P \Rightarrow Q) \land (Q \Rightarrow P).$$
\end{exo}

\begin{exo}[Calcul avec valeurs donn\'ees \corrige]{1}
On pose $P = \vrai$, $Q = \faux$, $R = \vrai$.
Calculer la valeur de v\'erit\'e de chaque formule.
\begin{multicols}{2}
\begin{enumerate}
  \item $P \land Q$
  \item $P \lor Q$
  \item $P \Rightarrow Q$
  \item $Q \Rightarrow P$
  \item $(P \lor Q) \Rightarrow R$
  \item $\lnot P \lor (Q \land R)$
  \item $P \Leftrightarrow R$
  \item $(P \land \lnot Q) \Rightarrow R$
\end{enumerate}
\end{multicols}
\end{exo}

\begin{exo}[Forme logique de r\'esultats math\'ematiques \corrige]{1}
R\'e\'ecrire chacun des r\'esultats suivants sous la forme
$\forall \ldots,\ P \Rightarrow Q$ (ou $P \Leftrightarrow Q$).
\begin{enumerate}
  \item Le carr\'e de tout entier impair est impair.
  \item Un triangle est rectangle si et seulement si le carr\'e de son
        plus grand c\^ot\'e \'egal la somme des carr\'es des deux autres.
  \item Si $n$ est divisible par $2$ et par $3$, alors $n$ est divisible par $6$.
  \item Tout nombre de la forme $4k+2$ est divisible par $2$.
\end{enumerate}
\end{exo}

\begin{exo}[Priorit\'e des connecteurs \corrige]{1}
Les connecteurs logiques ob\'eissent \`a des priorit\'es :
$\lnot$ (le plus prioritaire), puis $\land$, puis $\lor$, puis $\Rightarrow$,
puis $\Leftrightarrow$ (le moins prioritaire).

Mettre les parenth\`eses implicites dans les formules suivantes,
puis dresser leur table de v\'erit\'e.
\begin{enumerate}
  \item $\lnot P \land Q \lor R$
  \item $P \land Q \Rightarrow R$
  \item $P \lor Q \land \lnot R$
\end{enumerate}
\end{exo}

\subsection*{\diff{2}\quad Niveau 2 \textemdash\ Combinaisons et raisonnements}

\begin{exo}[\logiq\ De Morgan : simplifications \corrige]{2}
Simplifier en utilisant les lois de De Morgan\index{lois de De Morgan}.
\emph{Rappel :} $\lnot(P \land Q) \equiv \lnot P \lor \lnot Q$ ;
$\lnot(P \lor Q) \equiv \lnot P \land \lnot Q$.
\begin{multicols}{2}
\begin{enumerate}
  \item $\lnot(\lnot P \land Q)$
  \item $\lnot(P \lor \lnot Q)$
  \item $\lnot\bigl((P \Rightarrow Q) \land P\bigr)$
  \item $\lnot(\lnot P \lor \lnot Q)$
  \item $\lnot\bigl(\lnot(P \land Q)\bigr)$
  \item $\lnot(P \Leftrightarrow Q)$
\end{enumerate}
\end{multicols}
\end{exo}

\begin{exo}[Table \`a trois variables \corrige]{2}
Dresser la table de v\'erit\'e compl\`ete ($8$ lignes) de
$(P \Rightarrow Q) \land (Q \Rightarrow R)$ et de $P \Rightarrow R$.
Que constate-t-on~? Conclure sur la \textbf{transitivit\'e} de l'implication.
\end{exo}

\begin{exo}[\logiq\ Tautologies et contradictions \corrige]{2}
V\'erifier par table de v\'erit\'e\index{tautologie}\index{contradiction} :
\begin{enumerate}
  \item $P \lor \lnot P$ est une tautologie.
  \item $P \land \lnot P$ est une contradiction.
  \item $P \Rightarrow (Q \Rightarrow P)$ est une tautologie.
  \item $(P \Rightarrow Q) \Leftrightarrow (\lnot P \lor Q)$ est une tautologie.
\end{enumerate}
\end{exo}

\begin{exo}[Divisibilit\'e et implication \corrige]{2}
Soit $P$ : \og $n$ est divisible par $4$\fg\ et
$Q$ : \og $n$ est pair\fg.
\begin{enumerate}
  \item $P \Rightarrow Q$ est-elle vraie~? Et $Q \Rightarrow P$~?
  \item L'\'equivalence $P \Leftrightarrow Q$ est-elle vraie~?
  \item \'Ecrire et v\'erifier directement la contrapos\'ee de $P \Rightarrow Q$.
  \item G\'en\'eraliser : si $a \mid b$ et $b \mid c$, que peut-on dire de $a$ et $c$~?
\end{enumerate}
\end{exo}

\begin{exo}[Quantificateurs imbriqu\'es \corrige]{2}
Traduire en fran\c cais, d\'eterminer la valeur de v\'erit\'e,
puis \'ecrire la n\'egation.
\begin{enumerate}
  \item $\forall x \in \R,\; \exists y \in \R,\; y = x + 1$
  \item $\exists x \in \R,\; \forall y \in \R,\; y = x + 1$
  \item $\forall x \in \R,\; \exists y \in \R,\; x + y = 0$
  \item $\exists x \in \R,\; \forall y \in \R,\; x \cdot y = y$
\end{enumerate}
L'ordre des quantificateurs change-t-il le sens~?
\end{exo}

\begin{exo}[\logiq\ Raisonnement par contrapos\'ee \corrige]{2}
D\'emontrer par contrapos\'ee\index{raisonnement par contraposée}.
\begin{enumerate}
  \item Si $n^2$ est impair, alors $n$ est impair.
  \item Si $n^2$ est pair, alors $n$ est pair.
  \item Si $n$ est tel que $n^2 + n$ soit impair, alors $n$ est pair.
  \item Si $ab \ne 0$, alors $a \ne 0$ et $b \ne 0$.
\end{enumerate}
\end{exo}

\begin{exo}[G\'eom\'etrie et logique \corrige]{2}
Soient $P$ : \og le triangle $ABC$ est isoc\`ele en $A$\fg\ et
$Q$ : \og les angles $\hat B$ et $\hat C$ sont \'egaux\fg.
\begin{enumerate}
  \item Formuler $P \Rightarrow Q$ en fran\c cais. Est-ce vrai~?
  \item Formuler $Q \Rightarrow P$. Conclure sur $P \Leftrightarrow Q$.
  \item D\'emontrer la contrapos\'ee de $P \Rightarrow Q$.
  \item \'Enoncer $\lnot P$ et $\lnot Q$.
\end{enumerate}
\end{exo}

\begin{exo}[N\'egation de propositions complexes \corrige]{2}
\'Ecrire la n\'egation des propositions suivantes en langage naturel
\emph{et} avec des symboles.
\begin{enumerate}
  \item \og Tout entier pair est divisible par $4$ ou par $6$.\fg
  \item \og Il existe un entier $n$ tel que $n^2 < n < n^3$.\fg
  \item \og Pour tout r\'eel $x > 0$, il existe $y > 0$ tel que $y < x$.\fg
  \item \og Tout nombre premier est impair ou \'egal \`a $2$.\fg
\end{enumerate}
\end{exo}

\begin{exo}[Erreurs classiques de raisonnement \corrige]{2}
Identifier l'erreur logique\index{erreur de raisonnement} dans chaque argument.
\begin{enumerate}
  \item \og Tous les chats sont des animaux. Tous les chiens sont des animaux.
        Donc tous les chats sont des chiens.\fg
  \item \og Si $x^2 = 4$, alors $x = 2$. Or $(-2)^2 = 4$. Donc $-2 = 2$.\fg
  \item \og J'ai v\'erifi\'e pour $n=1,\ldots,5$ que $n^2+n+41$ est premier.
        Donc c'est toujours premier.\fg\ (Tester $n=40$.)
  \item Preuve que $1 = 2$ : on pose $a=b=1$.
        Alors $a^2 = ab$, donc $a^2 - b^2 = ab - b^2$,
        donc $(a-b)(a+b) = b(a-b)$,
        donc $a+b = b$, donc $1+1=1$, donc $2=1$.
        O\`u est l'erreur~?
\end{enumerate}
\end{exo}

\begin{exo}[Double implication \corrige]{2}
D\'emontrer l'\'equivalence en montrant les deux implications s\'epar\'ement~:
$$n \text{ est pair} \;\Leftrightarrow\; n^2 \text{ est pair}.$$
\end{exo}

\begin{exo}[Conditions n\'ecessaires et suffisantes \corrige]{2}
\index{condition nécessaire}\index{condition suffisante}
Pour chaque paire, pr\'eciser si la premi\`ere est une condition
n\'ecessaire, suffisante, ou les deux (\textit{c.n.s.}) pour la seconde.
\begin{enumerate}
  \item $n$ divisible par $6$ ~;~ $n$ divisible par $2$.
  \item $x = 0$ ~;~ $xy = 0$.
  \item $ABCD$ est un carr\'e ~;~ $ABCD$ est un rectangle.
  \item $\Delta = 0$ ~;~ le trin\^{o}me a une racine double.
  \item $n$ impair ~;~ $n^2$ impair.
\end{enumerate}
\end{exo}

\begin{exo}[D\'emonstration directe \corrige]{2}
D\'emontrer les assertions suivantes par un raisonnement direct.
\begin{enumerate}
  \item Si $n$ est impair, alors $n^3$ est impair.
  \item Si $n$ est impair, alors $n^2 + n$ est pair.
  \item Pour tout $a, b \in \Z$ : si $a$ et $b$ sont impairs,
        alors $a + b$ est pair.
  \item Pour tous $a, b, c \in \Z$ : si $a \mid b$ et $b \mid c$, alors $a \mid c$.
\end{enumerate}
\end{exo}

\begin{exo}[Trouver l'erreur dans une preuve \corrige]{2}
Chacune des \og d\'emonstrations\fg\ suivantes contient une erreur.
La trouver, l'expliquer et corriger la preuve.
\begin{enumerate}
  \item \textbf{Affirmation :} Tout entier v\'erifie $n^2 \geq n$.\\
        \textbf{Preuve :} On divise par $n$ pour obtenir $n \geq 1$. \'Evident.
  \item \textbf{Affirmation :} $\sqrt{a^2+b^2} = a + b$ pour $a,b \geq 0$.\\
        \textbf{Preuve :} En \'elevant au carr\'e les deux membres~:
        $a^2+b^2 = (a+b)^2 = a^2+2ab+b^2$, ce qui simplifie en $0=2ab$.
        Cela est vrai pour $a=0$ ou $b=0$. Donc c'est prouv\'e.
\end{enumerate}
\end{exo}

\subsection*{\diff{3}\quad Niveau 3 \textemdash\ D\'emonstrations et profondeur}

\begin{exo}[\logiq\ Irrationnalit\'e de $\sqrt{2}$ \corrige]{3}
D\'emontrer par l'absurde\index{raisonnement par l'absurde} que $\sqrt{2}$
est irrationnel.
\emph{On supposera $\sqrt{2} = p/q$ avec $\pgcd(p,q)=1$.}

Montrer ensuite que pour tout entier $k\geq 2$, si $k$ n'est pas un
carr\'e parfait, alors $\sqrt{k}$ est irrationnel.
\end{exo}

\begin{exo}[\logiq\ Irrationnalit\'e de $\sqrt{3}$ \corrige]{3}
D\'emontrer par l'absurde que $\sqrt{3}$ est irrationnel.
\emph{Indication : montrer d'abord que si $3 \mid n^2$, alors $3 \mid n$.}
\end{exo}

\begin{exo}[Suite et logique \corrige]{3}
Soit $P(n)$ : \og $n^2 - n$ est pair.\fg
\begin{enumerate}
  \item V\'erifier $P(n)$ pour $n = 0, 1, 2, 3, 4$.
  \item D\'emontrer $\forall n \in \N,\ P(n)$ en \'ecrivant $n^2-n = n(n-1)$.
  \item En d\'eduire : $\forall n \in \Z$, le carr\'e de $n$ est de la forme
        $4k$ ou $4k+1$ pour un certain entier $k$.
\end{enumerate}
\end{exo}

\begin{exo}[\logiq\ Infinit\'e des nombres premiers \corrige]{3}
D\'emontrer par l'absurde qu'il existe une infinit\'e de nombres premiers.
\index{nombre premier!infinité}

\emph{Indication (Euclide) : supposer qu'il n'en existe qu'un nombre fini
$p_1, p_2, \ldots, p_k$ et consid\'erer $N = p_1 p_2 \cdots p_k + 1$.}
\end{exo}

\begin{exo}[Combinaisons logiques \corrige]{3}
D\'emontrer (par table de v\'erit\'e ou directement) :
\begin{enumerate}
  \item $P \Rightarrow (Q \Rightarrow R)
        \;\equiv\; (P \land Q) \Rightarrow R$\quad\emph{(exportation)}
  \item $(P \Rightarrow Q) \land (P \Rightarrow R)
        \;\equiv\; P \Rightarrow (Q \land R)$
  \item $P \lor (Q \land R)
        \;\equiv\; (P \lor Q) \land (P \lor R)$
        \quad\emph{(distributivit\'e)}
\end{enumerate}
\end{exo}

\begin{exo}[Raisonner sur les r\'eels \corrige]{3}
\begin{enumerate}
  \item D\'emontrer $\forall x \in \R,\ x^2 \geq 0$.
        En d\'eduire $\forall x \in \R,\ x^2 + 1 > 0$.
  \item D\'emontrer par l'absurde qu'il n'existe pas de plus petit
        r\'eel strictement positif.
  \item D\'emontrer $\forall x, y \in \R,\ xy \leq \dfrac{x^2+y^2}{2}$.
\end{enumerate}
\end{exo}

\begin{exo}[Nombres rationnels et puissances \corrige]{3}
\begin{enumerate}
  \item Montrer que $2^n \ne 3^m$ pour tous entiers $n, m \geq 1$.
        \emph{(Indication~: $2^n$ est pair, $3^m$ est impair.)}
  \item En d\'eduire que la suite $(2^n/3^n)_{n\geq1}$ ne contient
        jamais un entier.
  \item Montrer que $(2/3)^n \to 0$ en v\'erifiant num\'eriquement
        pour $n = 1, 5, 10, 20$.
        L'expression $(2/3)^n$ peut-elle \^etre nulle~?
\end{enumerate}
\end{exo}

\begin{exo}[R\'eciproque du th\'eor\`eme de Pythagore \corrige]{3}
\textbf{Th\'eor\`eme de Pythagore :} Si un triangle $ABC$ est rectangle en $C$,
alors $AB^2 = AC^2 + BC^2$.

\begin{enumerate}
  \item \'Enoncer la contrapos\'ee du th\'eor\`eme de Pythagore.
  \item \'Enoncer la r\'eciproque : si $AB^2 = AC^2 + BC^2$, alors le
        triangle est rectangle en $C$.
  \item En admettant la r\'eciproque, d\'eterminer si le triangle de c\^ot\'es
        $5$, $12$, $13$ est rectangle, puis celui de c\^ot\'es $6$, $8$, $10$.
  \item La contrapos\'ee de la r\'eciproque est-elle vraie~? La formuler.
\end{enumerate}
\end{exo}

%─────────────────────────────────────────────────────────────────
\subsection*{Probl\`emes}
%─────────────────────────────────────────────────────────────────

\begin{probleme}[Circuit \`a trois interrupteurs \corrigeP]{1}
On code par \vrai{} un interrupteur ferm\'e et par \faux{} un interrupteur ouvert.
Le courant passe exactement lorsque la proposition suivante est vraie~:
$$\Phi(A,B,C) = \bigl(A \land (B \lor C)\bigr) \lor (\lnot A \land B \land C).$$
\begin{enumerate}
  \item Recopier et compl\'eter la table de v\'erit\'e ($8$ configurations possibles).
  \item Relever les configurations pour lesquelles le courant passe.
  \item Montrer que $\Phi \equiv (A\land B)\lor(A\land C)\lor(B\land C)$.
  \item Le courant passe-t-il exactement lorsqu'au moins deux interrupteurs sont ferm\'es~?
        Cette r\`egle est appel\'ee \emph{fonction majorit\'e}.
\end{enumerate}
\end{probleme}

\begin{probleme}[Chevaliers et coquins \corrigeP]{1}
Sur l'\^ile des logiciens, les chevaliers disent toujours vrai
et les coquins mentent toujours. Vous rencontrez $X$, $Y$, $Z$.
\begin{itemize}
  \item $X$ d\'eclare : \og $Y$ est un coquin.\fg
  \item $Y$ d\'eclare : \og $X$ et $Z$ sont du m\^eme type.\fg
\end{itemize}
\begin{enumerate}
  \item Poser $P$ : \og $X$ est chevalier\fg\ et $Q$ : \og $Y$ est chevalier\fg.
        Traduire les d\'eclarations en propositions logiques.
  \item Analyser les quatre cas $(P,Q)\in\{\vrai,\faux\}^2$.
  \item En d\'eduire la nature de $X$, $Y$, $Z$.
\end{enumerate}
\end{probleme}

\begin{probleme}[Portes logiques universelles \corrigeP]{2}
Une porte \textsc{nand} calcule $\lnot(P \land Q)$ ;
une porte \textsc{nor} calcule $\lnot(P \lor Q)$.
\begin{enumerate}
  \item Exprimer $\lnot P$ en utilisant uniquement des portes \textsc{nand}.
  \item Exprimer $P \land Q$ en utilisant uniquement des portes \textsc{nand}.
  \item Exprimer $P \lor Q$ en utilisant uniquement des portes \textsc{nand}.
  \item Conclure : les portes \textsc{nand} suffisent \`a calculer toute
        fonction logique. On dit qu'elles forment un syst\`eme \emph{fonctionnellement complet}.
\end{enumerate}
\end{probleme}

\begin{probleme}[Enigme logique : qui a fait quoi~? \corrigeP]{1}
Trois suspects $A$, $B$, $C$ sont interrog\'es apr\`es un vol.
Chacun fait une d\'eclaration :
\begin{itemize}
  \item $A$ : \og Ce n'est pas moi.\fg
  \item $B$ : \og C'est $A$.\fg
  \item $C$ : \og Ce n'est pas $A$.\fg
\end{itemize}
On sait qu'exactement une d\'eclaration est vraie.
D\'eterminer qui a commis le vol.
\end{probleme}

\begin{probleme}[Le berger, le loup, la ch\`evre et le chou \corrigeP]{1}
Un berger doit traverser une rivi\`ere avec un loup, une ch\`evre et un chou
en un seul voyage par trajet. Sa barque ne peut transporter qu'un seul passager.
Sans surveillance, le loup mange la ch\`evre, et la ch\`evre mange le chou.
\begin{enumerate}
  \item Mod\'eliser l'\'etat du probl\`eme par un quadruplet de propositions
        $(B, L, C, X)$ ($\vrai$ = c\^ot\'e d'arriv\'ee, $\faux$ = c\^ot\'e de d\'epart).
  \item Lister tous les \'etats interdits.
  \item Trouver une solution en $7$ traversées.
  \item Y a-t-il plusieurs solutions~?
\end{enumerate}
\end{probleme}

\begin{probleme}[Parit\'e des puissances \corrigeP]{2}
\begin{enumerate}
  \item D\'emontrer que pour tout entier $n$,
        $n$ et $n^3$ ont la m\^eme parit\'e.
  \item D\'emontrer que si $n$ est impair, alors $n^2 = 8k+1$
        pour un certain entier $k$.
        \emph{(Indication~: $n = 2m+1$, d\'evelopper $n^2$.)}
  \item En d\'eduire que si la somme $a^2 + b^2$ est divisible par $4$,
        alors $a$ et $b$ sont tous les deux pairs.
\end{enumerate}
\end{probleme}

\begin{probleme}[Divisibilit\'e et parit\'e des carr\'es \corrigeP]{2}
\begin{enumerate}
  \item D\'emontrer que le carr\'e de tout entier est de la forme $4k$
        ou $4k+1$ pour un certain entier $k$.
        \emph{(Distinguer le cas $n$ pair et $n$ impair.)}
  \item En d\'eduire que la somme $a^2 + b^2$ ne peut jamais \^etre
        de la forme $4k+3$.
        Donc une somme de deux carr\'es ne peut pas \^etre \'egale \`a
        $3$, $7$, $11$, $15$, \ldots
  \item V\'erifier ces faits pour $a^2+b^2 \in \{1,2,3,4,5,6,7,8\}$~:
        lesquelles de ces valeurs sont effectivement impossibles~?
  \item Quels entiers de $\{1,\ldots,20\}$ sont sommes de deux carr\'es~?
        Dresser la liste.
\end{enumerate}
\end{probleme}

\begin{probleme}[Logique dans les preuves g\'eom\'etriques \corrigeP]{2}
On consid\`ere les propositions~:
$P$~: \og le quadrilat\`ere $ABCD$ est un carr\'e\fg~;
$Q$~: \og $ABCD$ est un rectangle\fg~;
$R$~: \og $ABCD$ est un losange\fg~;
$S$~: \og $ABCD$ est un parall\'elogramme\fg.
\begin{enumerate}
  \item Pour chaque paire, pr\'eciser si on a $\Rightarrow$, $\Leftarrow$
        ou $\Leftrightarrow$~: $(P,Q)$~; $(P,R)$~; $(Q,S)$~; $(R,S)$~; $(P,S)$.
  \item Formuler la contrapos\'ee de \og\(P \Rightarrow Q\)\fg{}
        et de \og\(Q \Rightarrow S\)\fg{}.
  \item D\'emontrer par contrapos\'ee~:
        \og Si $ABCD$ n'est pas un parall\'elogramme, alors ce n'est pas un carr\'e.\fg
  \item \'Ecrire la n\'egation de~:
        \og Tout carr\'e est \`a la fois un rectangle et un losange.\fg{}
        Puis montrer que cette n\'egation est fausse.
\end{enumerate}
\end{probleme}

\begin{probleme}[Sudoku $4\times 4$ par d\'eduction logique \corrigeP]{1}
La grille $4\times 4$ suivante doit \^etre compl\'et\'ee avec les chiffres $1$, $2$, $3$, $4$
de fa\c con que chaque ligne, colonne et bloc $2\times 2$ contienne chaque chiffre exactement une fois.
\[
\begin{array}{|cc|cc|}
\hline
  & 2 & 3 &   \\
3 &   &   & 2 \\
\hline
  & 3 & 4 &   \\
4 &   &   & 3 \\
\hline
\end{array}
\]
\begin{enumerate}
  \item Expliquer chaque \'etape de votre raisonnement sous forme d'implication
        (ex. : \og La case $(1,1)$ ne peut pas contenir $3$ car...\fg).
  \item Combien d'implications avez-vous utilis\'ees~?
  \item Y a-t-il une autre solution~?
\end{enumerate}
\end{probleme}

\begin{probleme}[Sommes et formules alg\'ebriques \corrigeP]{2}
\begin{enumerate}
  \item V\'erifier pour $n = 1, 2, 3, 4, 5$ que~:
        $1 + 2 + \cdots + n = \dfrac{n(n+1)}{2}.$
  \item V\'erifier la formule en calculant la somme des $n$ premiers entiers
        de deux fa\c cons~: directement, et via la formule.
        Sont-elles \'egales~?
  \item \'Ecrire un algorithme qui calcule $S = 1+2+\cdots+n$ par une boucle,
        puis v\'erifie si $S = n(n+1)/2$. Tester pour $n = 10, 100$.
  \item En observant le tableau de valeurs, conjecturer une formule
        pour $1 + 3 + 5 + \cdots + (2n-1)$ (somme des $n$ premiers impairs).
        V\'erifier pour $n = 1, 2, 3, 4, 5$.
\end{enumerate}
\end{probleme}

\begin{probleme}[Codage et logique~: parit\'e \corrigeP]{2}
Un bit de parit\'e est ajout\'e \`a un message binaire pour d\'etecter les erreurs~:
le nombre total de $1$ dans le message \emph{plus} le bit de parit\'e doit \^etre pair.
\begin{enumerate}
  \item Le message \texttt{1011} a-t-il un nombre pair ou impair de $1$~?
        Quel bit de parit\'e faut-il ajouter~?
  \item Si on re\c coit \texttt{10110} (message~+~parit\'e), y a-t-il une erreur~?
  \item Si on re\c coit \texttt{10100}, y a-t-il une erreur~?
  \item Ce syst\`eme peut-il d\'etecter \emph{deux} erreurs simultan\'ees~?
        Expliquer.
\end{enumerate}
\end{probleme}

\begin{probleme}[D\'eduction logique~: tournoi de foot \corrigeP]{1}
Dans un tournoi \`a la ronde, chaque \'equipe joue contre chacune des autres
exactement une fois. Il y a $4$ \'equipes~: $A$, $B$, $C$, $D$.
Les r\`egles~: une victoire donne $3$ points, un nul $1$ point, une d\'efaite $0$.
On sait que~:
\begin{itemize}
  \item $A$ a gagn\'e ses $3$ matchs.
  \item $B$ a obtenu exactement $6$ points.
  \item $C$ et $D$ ont chacune le m\^eme nombre de points.
\end{itemize}
\begin{enumerate}
  \item Combien de matchs ont \'et\'e jou\'es en tout~?
  \item Combien de points $A$ a-t-elle~?
  \item D\'eterminer les points de $B$, $C$ et $D$.
        \emph{(Attention~: une victoire distribue $3$ points, mais un match nul
        n'en distribue que $2$.)}
  \item Reconstituer le classement final.
\end{enumerate}
\end{probleme}

%─────────────────────────────────────────────────────────────────
\subsection*{D\'efis}
%─────────────────────────────────────────────────────────────────

\begin{challenge}[Le paradoxe du menteur]
\index{paradoxe du menteur}
Consid\'erons la phrase~: \og Cette phrase est fausse.\fg
\begin{enumerate}
  \item Si on suppose qu'elle est \textbf{vraie}, qu'en d\'eduit-on~?
        Si on suppose qu'elle est \textbf{fausse}, qu'en d\'eduit-on~?
  \item Pourquoi cette phrase n'est-elle pas une proposition au sens
        du cours~?
  \item Inventer deux autres phrases auto-r\'ef\'erentes.
        Sont-elles des propositions~? Lesquelles~?
        \begin{itemize}
          \item \og Cette phrase contient cinq mots.\fg
          \item \og Cette phrase est vraie.\fg
        \end{itemize}
  \item Le \emph{paradoxe de Russell}~: consid\'erons \og l'ensemble de tous
        les ensembles qui ne se contiennent pas eux-m\^emes\fg.
        Cet ensemble se contient-il lui-m\^eme~?
        En quoi retrouve-t-on la m\^eme structure que le paradoxe du menteur~?
\end{enumerate}
\end{challenge}

\begin{challenge}[Construire sa propre table de v\'erit\'e]
Un \emph{ou exclusif} (not\'e $P \oplus Q$) est vrai quand exactement
l'une des deux propositions est vraie (mais pas les deux).
\begin{enumerate}
  \item Dresser la table de v\'erit\'e de $P \oplus Q$.
  \item V\'erifier que $P \oplus Q \equiv (P \lor Q) \land \lnot(P \land Q)$.
  \item Exprimer $P \oplus Q$ uniquement avec $\Rightarrow$ et $\lnot$.
        \emph{(Indication~: partir de $(P\lor Q)\land\lnot(P\land Q)$
        et utiliser $(P\Rightarrow Q)\equiv(\lnot P\lor Q)$.)}
  \item V\'erifier que $P \oplus P \equiv \faux$ et $P \oplus \lnot P \equiv \vrai$
        pour toutes valeurs de $P$.
  \item Application~: le chiffrement XOR. Si un message est cod\'e
        par le bit $M$ et une cl\'e secrète $K$, le chiffr\'e est $C = M \oplus K$.
        Montrer que $C \oplus K = M$ (le d\'echiffrement est la m\^eme op\'eration).
\end{enumerate}
\end{challenge}

\begin{challenge}[Enigmes logiques~: chevaliers et coquins]
\index{chevaliers et coquins}
Sur l'\^ile des logiciens, les \emph{chevaliers} disent toujours vrai
et les \emph{coquins} mentent toujours.
\begin{enumerate}
  \item Vous rencontrez $A$ et $B$.
        $A$ dit~: \og Nous sommes tous les deux des coquins.\fg
        D\'eterminer la nature de $A$ et de $B$.
        \emph{(Raisonner par cas~: si $A$ est chevalier\ldots)}
  \item Vous rencontrez $C$, $D$, $E$.
        $C$ dit~: \og Au moins l'un d'entre nous est un coquin.\fg
        $D$ dit~: \og$C$ est un coquin.\fg
        $E$ ne dit rien.
        Quelles sont les natures possibles de $C$, $D$, $E$~?
  \item Inventer une \'enigme du m\^eme type avec $3$ personnes
        et \textbf{exactement une} solution.
  \item Formaliser la question~(1) \`a l'aide de propositions logiques
        ($P$~: \og$A$ est chevalier\fg, $Q$~: \og$B$ est chevalier\fg)
        et r\'esoudre par table de v\'erit\'e.
\end{enumerate}
\end{challenge}

\begin{challenge}[Bijections et taille d'ensembles]
\index{bijection}\index{Cantor, Georg}
Deux ensembles ont \og autant\fg\ d'\'el\'ements si l'on peut associer
\`a chaque \'el\'ement de l'un exactement un \'el\'ement de l'autre
(c'est une \emph{bijection}).
\begin{enumerate}
  \item Les ensembles $\{1,2,3\}$ et $\{a,b,c\}$ ont-ils autant
        d'\'el\'ements~? Exhiber une bijection.
  \item L'ensemble $\N = \{0,1,2,3,\ldots\}$ et
        $\Z = \{\ldots,-2,-1,0,1,2,\ldots\}$
        ont-ils autant d'\'el\'ements~?
        \emph{Tentative de bijection~:}
        $f(0)=0,\; f(1)=1,\; f(2)=-1,\; f(3)=2,\; f(4)=-2,\ldots$
        \'Ecrire la formule explicite de $f(n)$ selon que $n$ est pair ou impair.
  \item Les entiers pairs $\{0,2,4,6,\ldots\}$ ont-ils autant d'\'el\'ements
        que $\N$ tout entier~?
        C'est le \emph{paradoxe de Galileo}~: une partie peut \^etre
        en bijection avec le tout~!
  \item Cantor ($1845$--$1918$) a montr\'e que $\R$ est \og plus grand\fg\ que $\N$~:
        il n'existe aucune bijection $\N\to\R$.
        Sans d\'emontrer cela, expliquer pourquoi il est impossible de
        \og lister\fg\ tous les r\'eels de $[0,1]$.
\end{enumerate}
\end{challenge}

\begin{probleme}[\logicoalgo\ Logique et calcul alg\'ebrique \corrigeP]{2}
Les propositions suivantes portent sur des r\'eels $x$ et $y$.
\begin{enumerate}
  \item \'Enoncer la n\'egation de~:
        \og Pour tout r\'eel $x$, $x^2 \geq 0$.\fg
        Puis v\'erifier que la proposition originale est vraie.
  \item \'Enoncer la contrapos\'ee de~:
        \og Si $x^2 = 0$ alors $x = 0$.\fg
        D\'emontrer la contrapos\'ee.
  \item \'Enoncer sous forme logique~:
        \og $ax^2+bx+c=0$ a deux racines distinctes si et seulement si
        $\Delta > 0$.\fg
        Quelle est la r\'eciproque~? Est-elle vraie~?
  \item \'Enoncer la n\'egation de \og $\exists x\in\R,\ x^2+1=0$\fg\ et
        v\'erifier que la n\'egation est vraie.
\end{enumerate}
\end{probleme}

\begin{probleme}[\logicoalgo\ Algorithme et divisibilit\'e \corrigeP]{2}
On cherche tous les entiers $n$ entre $1$ et $100$ divisibles par $3$ mais pas par $9$.

\begin{algorithm}[H]
\caption{Divisible par 3 mais pas par 9}
\KwIn{rien (on examine $n = 1, \ldots, 100$)}
\Pour{$n \leftarrow 1$ \KwTo $100$}{
  \Si{$n \mod 3 = 0$ \textbf{et} $n \mod 9 \ne 0$}{
    \Afficher $n$\;
  }
}
\end{algorithm}

\begin{enumerate}
  \item La condition \og$n\mod3=0$ et $n\mod9\ne0$\fg\ est-elle
        \'equivalente \`a \og$3\mid n$ et $9\nmid n$\fg~? Justifier.
  \item \'Enoncer la n\'egation de cette condition.
  \item Ex\'ecuter l'algorithme pour $n\in\{3,6,9,12,18,21\}$.
  \item Combien d'entiers entre $1$ et $100$ v\'erifient cette condition~?
        Donner la r\'eponse avec un argument logique sans tout \'enum\'erer.
\end{enumerate}
\end{probleme}

\begin{challenge}[Raisonnement par l'absurde~: deux preuves classiques]
\begin{enumerate}
  \item \textbf{Infinit\'e des multiples de $7$.}
        On suppose par l'absurde qu'il n'existe qu'un nombre fini de
        multiples de $7$~: $7, 14, 21, \ldots, 7k$ (le plus grand).
        Montrer que $7(k+1)$ est aussi un multiple de $7$, et conclure.
  \item \textbf{Il n'y a pas de plus grand entier.}
        Supposons qu'il existe un plus grand entier $N$.
        D\'eriver une contradiction.
  \item \textbf{$\sqrt{5}$ est irrationnel.}
        Adapter la preuve de $\sqrt{2}$ (exercice~\hyperref[exo:1.29]{1.29})~:
        si $\sqrt{5}=p/q$ irr\'eductible, montrer que $5\mid p$
        puis $5\mid q$.
  \item \textbf{Il n'existe pas deux r\'eels cons\'ecutifs.}
        Montrer que pour tous $a < b$ r\'eels, il existe toujours
        un r\'eel $c$ v\'erifiant $a < c < b$.
        \emph{(Indication~: prendre $c = (a+b)/2$.)}
\end{enumerate}
\end{challenge}

\begin{challenge}[Construire un syst\`eme de r\`egles logiques]
La r\`egle du \emph{modus ponens}\index{modus ponens} dit~:
si l'on sait que $P$ est vraie, et que $P \Rightarrow Q$ est vraie,
alors on peut conclure que $Q$ est vraie.
\begin{enumerate}
  \item Appliquer le modus ponens dans les situations suivantes~:
        \begin{itemize}
          \item $P$~: \og Il pleut.\fg\quad
                $P\Rightarrow Q$~: \og S'il pleut, le sol est mouill\'e.\fg\\
                Conclure~?
          \item $P$~: \og $n$ est divisible par $6$.\fg\quad
                $P\Rightarrow Q$~: \og Si $n$ est divisible par $6$, il est
                divisible par $2$.\fg\quad Conclure~?
        \end{itemize}
  \item Cha\^ine de raisonnement~: on dispose de~:
        (1)~$P$, \quad (2)~$P\Rightarrow Q$, \quad (3)~$Q\Rightarrow R$.
        D\'eduire $R$ en appliquant deux fois le modus ponens.
  \item Montrer que le syll\-o\-gisme (si tous les $A$ sont des $B$,
        et si tous les $B$ sont des $C$, alors tous les $A$ sont des $C$)
        est une application du modus ponens avec des quantificateurs.
  \item Construire soi-m\^eme un enchanement de $3$ implications vraies
        m\^enant d'une hypoth\`ese de d\'epart \`a une conclusion non triviale.
\end{enumerate}
\end{challenge}

%─────────────────────────────────────────────────────────────────
\begin{bilan}
\begin{itemize}
  \item \textbf{Proposition} : \'enonc\'e vrai ou faux
        (pas un pr\'edicat, ni un ordre).
  \item \textbf{Connecteurs} :
        $\lnot$ (non), $\land$ (et), $\lor$ (ou inclusif),
        $\Rightarrow$ (implique), $\Leftrightarrow$ (\'equivaut).
  \item \textbf{Implication} : fausse \emph{uniquement} quand
        $P=\vrai$ et $Q=\faux$.
  \item \textbf{Contrapos\'ee} : $P\Rightarrow Q\equiv\lnot Q\Rightarrow\lnot P$. \\
        La r\'eciproque $Q\Rightarrow P$ n'est pas \'equivalente.
  \item \textbf{Quantificateurs} : $\forall$ et $\exists$.
        N\'egation : $\lnot\forall \equiv \exists\lnot$ ;
        $\lnot\exists \equiv \forall\lnot$.
  \item \textbf{De Morgan} :
        $\lnot(P\land Q)\equiv\lnot P\lor\lnot Q$ ;
        $\lnot(P\lor Q)\equiv\lnot P\land\lnot Q$.
  \item \textbf{Raisonnements} : direct, contrapos\'ee, absurde.
        Un seul contre-exemple r\'efute un $\forall$.
\end{itemize}
\end{bilan}

%% ════════════════════════════════════════════════════════════════
